\documentclass{article}
\usepackage[preprint]{neurips_2026}
\usepackage[utf8]{inputenc}
\usepackage[T1]{fontenc}
\usepackage[spanish,es-tabla]{babel}
\usepackage{amsmath,amssymb,booktabs,graphicx,microtype,url}
\usepackage{hyperref}
\newcommand{\pending}[1]{\paragraph{Pendiente de evidencia.} #1}
\author{Grupo de cinco estudiantes: completar nombres y correos\\
Universidad Nacional de Ingeniería\\Lima, Perú}

\title{Clasificación multietiqueta de anomalías de rodilla con CNN 2.5D y supervisión débil multilingüe}
\begin{document}
\maketitle
\begin{abstract}
Se ejecutó una comparación factorial de ocho variantes de clasificación de doce hallazgos con CNN propia y ResNet-50, BCE/focal y dos tratamientos del ruido textual. Los 600 estudios incluyen 58 expertos reservados para evaluación. El AUC macro oficial quedó definido para las doce etiquetas. El baseline previamente fijado obtuvo AUC macro de doce etiquetas 0.578 sobre 12 etiquetas y F1 macro 0.361. Se auditaron predicciones y checkpoints y se generaron explicaciones para tres hallazgos. La muestra pequeña, el ruido lingüístico y una única semilla limitan las conclusiones.
\end{abstract}

\section{Motivación y trabajos relacionados}
Las imágenes y los reportes de RSNA constituyen fuentes complementarias de supervisión~\cite{rsna}.
Una etiqueta ausente en el reporte no implica que un radiólogo haya descartado el hallazgo.
El objetivo es evaluar cuánto afecta esta diferencia a la clasificación, manteniendo independiente
la selección del modelo respecto de los estudios expertos.

ResNet introduce conexiones residuales para facilitar la optimización de redes profundas~\cite{resnet}.
Se utiliza su implementación oficial como backbone preentrenado y se implementa una cabeza multietiqueta propia.
Focal loss reduce la contribución de ejemplos fáciles~\cite{focal}; aquí se adapta a doce salidas binarias,
con pesos por frecuencia calculados únicamente en entrenamiento. La supervisión débil requiere distinguir
entre reglas, abstenciones y etiquetas verificadas~\cite{snorkel}; este trabajo no implementa el modelo generativo
de Snorkel. Grad-CAM proporciona una visualización de la sensibilidad espacial~\cite{gradcam},
sin establecer que las regiones activadas tengan significado clínico causal.

\section{Datos y separación experimental}
Se validaron 4407 estudios y 24371 descriptores. La muestra efectiva tiene 434 estudios de entrenamiento, 108 de validación débil y 58 expertos. Se comprobaron hashes, identificadores y decodificación de 22954 DICOM. La caché coincide byte a byte con sus originales. Planos: \{'Sagittal': 600\}. Grupos de escritura de la muestra: \{'Latin/mixed': 532, 'Greek': 65, 'Cyrillic': 3\}.

El esquema observado en los CSV contiene 4\,407 estudios y 58 estudios con doce etiquetas binarias completas.
Se seleccionan hasta 600 estudios: todos los expertos y una muestra de reportes compatibles con el léxico.
La selección usa una puntuación de rareza multietiqueta y semilla 42; no garantiza positivos para etiquetas
sin evidencia textual. Las particiones se realizan por StudyInstanceUID; ningún corte de un estudio
se distribuye entre entrenamiento, validación y evaluación.

Los estudios no expertos se dividen en entrenamiento y validación débil, aproximadamente 80/20.
La versión final asigna cuotas de positivos por etiqueta a validación, conservando evidencia positiva
en entrenamiento. Una primera ejecución utilizó una asignación por rareza que dejó Lateral OA
sin positivos débiles en entrenamiento (0 frente a 14 en validación). Se conservó esa ejecución,
se corrigió la partición usando exclusivamente etiquetas débiles y se repitió toda la matriz.
Los mismos 58 expertos se evaluaron en ambas ejecuciones; la cohorte final no constituye
una nueva prueba externa nunca consultada. Las cuotas no se ajustaron por resultados expertos.
Los expertos se utilizan exclusivamente para evaluación y para medir el acuerdo de las reglas fijadas previamente.
No se cambia el léxico, el umbral ni los hiperparámetros tras observar esa evaluación.
Los CSV no proporcionan un identificador de paciente transversal. La auditoría de los PatientID
disponibles en DICOM no detectó intersecciones entre particiones, pero su carácter anonimizado
no permite garantizar independencia de personas reales entre estudios.

Se prefiere una serie sagital sensible a líquido con supresión grasa. Cualquier alternativa queda registrada.
Los DICOM se ordenan por geometría y se verifican orientación, identificadores y discontinuidades.
La alternativa basada en InstanceNumber genera una advertencia.
Se aplica la transformación de modalidad, corrección MONOCHROME1 y escalado por percentiles 0.5/99.5 del volumen.
Se toman ocho grupos de tres cortes contiguos, redimensionados a $224\times224$.
Los tres cortes actúan como canales de entrada; no representan color RGB real.

\section{Extracción y tratamiento de ruido}
El léxico cubre términos en inglés, español, francés, alemán, neerlandés, portugués, italiano,
turco, griego, búlgaro, croata, chino y tailandés. Es una cobertura de vocabulario, no una demostración
de competencia lingüística equivalente. Las cláusulas, entidades vecinas y ventanas locales limitan el alcance
de las negaciones. Incertidumbre o afirmaciones contradictorias producen abstención.
Una afirmación explícita recibe peso 0.9; una etiqueta no mencionada en un reporte compatible recibe 0.2;
un reporte sin entidades reconocidas queda completamente enmascarado.
Estos valores son pesos heurísticos, no probabilidades calibradas.

La tasa de desacuerdo se calcula por etiqueta entre predicciones textuales comparables y etiquetas expertas,
junto con su cobertura. Se presenta por separado para todas las reglas y para confianza $\geq0.8$.
El grupo de escritura del texto sirve para auditar cobertura y no se interpreta como idioma o centro de origen.

\section{Arquitecturas y pérdidas}
La CNN propia tiene cuatro bloques Conv--GroupNorm--ReLU--MaxPool con 32, 64, 128 y 256 canales.
Tras el pooling global se aplica dropout 0.2 y una capa lineal de doce logits.
La implementación contiene 391\,980 parámetros entrenables.
La alternativa usa ResNet-50 con pesos ImageNet locales y ajusta todo el backbone.
Su encoder y cabeza suman 23\,532\,620 parámetros entrenables; se usan los pesos oficiales
ImageNet V2 de torchvision, verificados mediante su SHA-256.
En ambos casos se promedian los logits de los ocho grupos y se aplica sigmoid por etiqueta.

Con $m_{ij}$ como peso de supervisión y $y_{ij}$ como etiqueta observada:
\begin{equation}
\mathcal{L}_{BCE}=
\frac{\sum_{ij}m_{ij}[-y_{ij}\log p_{ij}-(1-y_{ij})\log(1-p_{ij})]}
{\max(1,\sum_{ij}m_{ij})}.
\end{equation}
Para focal loss se multiplica cada término por $\alpha_{t,j}(1-p_{t,ij})^2$.
Los pesos positivos se calculan como $1-\mathrm{prevalencia}_j$ y se limitan al intervalo $[0.1,0.9]$.
En la variante filtrada, la prevalencia se calcula después de excluir las etiquetas de confianza inferior a 0.8.
La alternativa de alta confianza conserva únicamente etiquetas con confianza $\geq0.8$;
descarta estudios totalmente sin supervisión de entrenamiento y validación.

\section{Protocolo experimental}
La matriz factorial comprende dos arquitecturas, dos pérdidas y dos estrategias de ruido: ocho experimentos.
Todos usan semilla 42, AdamW~\cite{adamw}, tasa $3\times10^{-4}$, weight decay $10^{-4}$,
batch 4, hasta 50 épocas y parada temprana con paciencia 10 sobre pérdida de validación débil.
Los checkpoints incluyen optimizador, AMP, generadores y estados aleatorios.
La repetibilidad numérica entre hardware y versiones deberá verificarse; una semilla no la garantiza.
La pérdida de cada época promedia por estudio las pérdidas normalizadas de sus minibatches;
la validación mantiene orden fijo. No se aplican aumentos de imagen en esta comparación base.

El umbral de clasificación 0.5 se fija antes de evaluar como baseline simétrico.
Se reportan precisión, recall, F1 y AUC por etiqueta.
El AUC macro oficial se informa solamente si las doce etiquetas tienen ambas clases;
un promedio parcial se identifica explícitamente y no sustituye la métrica oficial.
Precisión, recall y F1 usan la convención zero\_division=0; un cero por ausencia de predicciones
positivas se interpreta junto con el soporte, no como evidencia de rendimiento estable.
El análisis de prevalencia documenta las etiquetas poco frecuentes y la comparación focal/BCE mide
una estrategia de mitigación. La pequeña evaluación experta limita especialmente las estimaciones de hallazgos raros.

\section{Resultados e interpretabilidad}
\InputIfFileExists{tables.tex}{}{\pending{Ejecutar la matriz completa y generar las tablas desde métricas reales.}}
\InputIfFileExists{figures.tex}{}{\pending{Generar curvas, prevalencia y Grad-CAM de LCA, menisco medial y derrame.}}
En esta evaluación se dispone de AUC para 12 de los doce hallazgos. El AUC macro de doce etiquetas de las variantes varía entre 0.513 y 0.620; F1 macro varía entre 0.361 y 0.506. Estos intervalos describen la matriz ejecutada y no son intervalos de confianza. No se elige una variante para despliegue por estos resultados de evaluación. Las curvas muestran el criterio aplicado para elegir cada checkpoint; no se comparan los valores absolutos de BCE y focal como si fueran la misma escala.  Comparaciones pareadas por arquitectura y tratamiento del ruido: CNN/bce/pond.: diferencia focal menos BCE de AUC -0.005 y de F1 macro +0.135; CNN/bce/filtro: diferencia focal menos BCE de AUC +0.007 y de F1 macro -0.028; ResNet/bce/pond.: diferencia focal menos BCE de AUC +0.001 y de F1 macro +0.144; ResNet/bce/filtro: diferencia focal menos BCE de AUC -0.014 y de F1 macro +0.035. Con una semilla y una cohorte experta pequeña, estas diferencias no establecen superioridad general. La mitigación del desbalance debe evaluarse por hallazgo y soporte, incluyendo posibles pérdidas de precisión al elevar recall. El umbral 0,5 permanece fijo aun cuando una variante produce probabilidades mal calibradas. La distribución tras el filtro explica qué etiquetas conservan positivos y negativos explícitos; no se interpreta el filtrado como limpieza perfecta.

Medial OA y Lateral OA tienen respectivamente 35/434 (8.1\%) y 10/434 (2.3\%) positivos débiles
en entrenamiento. En ResNet con ponderación, cambiar BCE por focal eleva el recall experto de
Medial OA de 0.133 a 0.867 y F1 de 0.222 a 0.473, mientras la precisión baja de 0.667 a 0.325.
Para Lateral OA, recall pasa de 0 a 0.182 y F1 de 0 a 0.222, pero AUC disminuye de 0.683 a 0.644.
Esta mitigación modifica el compromiso al umbral fijo; no mejora uniformemente la discriminación
de hallazgos raros. Los soportes expertos de 15 y 11 positivos limitan la estabilidad de estas cifras.

El baseline cualitativo se fija como CNN propia, BCE y ponderación, evitando seleccionarlo por el conjunto experto.
Los mapas se calculan respecto del logit de cada hallazgo y muestran el grupo de cortes con mayor activación media.
En los tres casos mostrados, la probabilidad supera 0.5 y la etiqueta experta es positiva;
esto no hace que su localización sea correcta. La inspección gráfica del caso LCA muestra una
respuesta amplia y zonas periféricas, sin un foco que permita atribuir inequívocamente la señal al ligamento.
En menisco medial, la respuesta cubre de forma difusa el tejido visible y no delimita una región
meniscal selectiva. En derrame, los focos destacados son periféricos e incluyen fondo fuera del tejido;
tampoco ofrecen evidencia convincente de localización del hallazgo.
La heurística de máximo promedio seleccionó grupos extremos: 7 para LCA y 0 para menisco y derrame,
en índices de 0 a 7. Este criterio puede privilegiar extensión del mapa y fondo en lugar de una
estructura relevante; es una limitación del protocolo de visualización, no una localización clínica validada.
Los ocho mapas por caso se conservan en archivos NPZ para permitir revisar el volumen muestreado completo.
Estas observaciones son exploratorias: no hay regiones objetivo anotadas, métricas de localización ni
revisión radiológica profesional registrada. La resolución del último bloque de la CNN es $14\times14$;
interpolar a $224\times224$ suaviza el mapa sin añadir precisión anatómica. El color se normaliza por
grupo de cortes, por lo que su intensidad no compara causalmente los tres hallazgos. Una evaluación
posterior debería confrontar cada mapa con regiones marcadas por un especialista y añadir falsos positivos
y falsos negativos, sin cambiar esta evaluación final ni seleccionar ejemplos por apariencia.

\section{Discusión, sesgo y limitaciones}
La calidad de supervisión depende de la cobertura del léxico, abreviaturas y alcance de la negación.
La identificación de una escritura no equivale a determinar el idioma: el grupo latino reúne varias
lenguas y textos mixtos. Incluso cuando una entidad se reconoce, una construcción coordinada,
una negación distante, terminología local o un matiz de incertidumbre pueden producir errores.
Una traducción literal tampoco demuestra equivalencia de criterios clínicos. El desacuerdo de las reglas
debe leerse junto con la cobertura: una regla selectiva puede aparentar mejor acuerdo al abstenerse
en los casos difíciles. Los 58 expertos no permiten validar de forma sólida todas las combinaciones
de idioma, centro y hallazgo; los CSV no identifican de manera fiable los 16 centros ni su geografía.

La selección por rareza busca cobertura de positivos y modifica la distribución respecto al universo
de estudios. Las métricas no son estimaciones de prevalencia poblacional, y el filtrado adicional por
confianza puede concentrar ciertos estilos de reporte. Las imágenes asociadas a textos con menor
cobertura tienen menos supervisión, de modo que una desigualdad lingüística puede convertirse en
desigualdad de aprendizaje visual. La ponderación preserva evidencia de baja confianza, pero también
introduce falsos negativos al tratar la ausencia de mención como un negativo débil. El filtro elimina
esa evidencia y puede dejar hallazgos con muy pocos negativos explícitos.

El riesgo de despliegue consiste en aprender correlaciones de adquisición y de anotación. Una
normalización por volumen no elimina diferencias de equipos, secuencias o poblaciones. Son necesarias
etiquetas expertas adicionales estratificadas por idioma y centro, auditoría de negación por hablantes
competentes y evaluación externa con revisión profesional. El presente ejercicio no acredita un sistema
apto para uso clínico. Los mapas de activación ofrecen indicios espaciales y no sustituyen esa revisión.

Una normalización uniforme no elimina diferencias entre protocolos de adquisición.
Usar una sola serie y un muestreo 2.5D puede omitir hallazgos visibles en otros planos.
La validación débil comparte el mecanismo de ruido con entrenamiento y puede favorecer las mismas reglas erróneas.
El análisis de acuerdo es diagnóstico y no mide sensibilidad clínica en una población independiente.
Los reportes y las imágenes permanecen en el equipo local; no se publican con el código.

\section{Reproducibilidad y trabajo pendiente}
El código, la configuración congelada y los comandos se documentan en MRI\_RUN.md.
El repositorio del curso es \url{https://github.com/miguelsff/maestria-ia-uni-ciclo-2-redes-neuronales};
la ejecución verificada corresponde a la copia local conservada, pendiente de publicación de estos cambios.
El notebook se ejecuta leyendo los resultados persistidos, con los flags de entrenamiento desactivados.
La caché local de Docker contiene copias de las series elegidas verificadas por SHA-256;
los originales se mantienen montados en solo lectura. Las predicciones se recalculan durante la auditoría
de métricas y se verifica la selección del checkpoint contra las curvas registradas.
La prueba de reanudación verifica igualdad exacta de los parámetros con una ejecución continua,
en un fixture controlado; no demuestra igualdad entre versiones o equipos distintos.
\begin{table}[ht]
\centering\small
\caption{Tiempo de entrenamiento y evaluación, y máximo de memoria asignada por PyTorch; no equivalen al consumo total del dispositivo.}
\begin{tabular}{lrr}
\toprule
Variante & Segundos & VRAM tensor GiB \\
\midrule
CNN/bce/pond. & 115.5191 & 1.0720 \\
CNN/bce/filtro & 54.7096 & 1.0738 \\
CNN/focal/pond. & 70.3960 & 1.0730 \\
CNN/focal/filtro & 56.4852 & 1.0738 \\
ResNet/bce/pond. & 67.0683 & 1.7379 \\
ResNet/bce/filtro & 60.8994 & 1.7352 \\
ResNet/focal/pond. & 68.6879 & 1.7352 \\
ResNet/focal/filtro & 68.8144 & 1.7352 \\
\bottomrule
\end{tabular}
\end{table}

Las guías de habilidades procedimentales orientaron la preparación y la evaluación~\cite{skills};
no constituyen validación experimental. El presente PDF conserva estado de borrador hasta completar
los cinco autores y correos reales, revisión anatómica por una persona competente y revisión editorial del grupo.
El código local aún debe publicarse en el repositorio por decisión de sus autores.
\bibliographystyle{plain}
\bibliography{references}
\end{document}
